% Fragment: fig-sealed-laundering-fork.tex -- the laundering residual as a
% forking-path comparison.
%
% Reader question: how could a worker defeat the "committed input" gate
% honestly, and why can't the current checker catch it?
% Claim: computing t=f(s) and submitting t lets the honest gate release
% g(t) instead of g(s); the fork sits strictly before the one node
% (submit) the checker's model can see, so the checker's search never
% notices anything changed.
% Evidence: the boundary after Theorem sealed-noninterference ("There is a
% second residual...it does not cover a worker that launders...the
% model's submit carries no payload") and Exercise sealed-launder's
% solution (a compute(f) transition, and a gate that then releases
% g(t)), website-v2/public/whitepaper/sealed-harbor.tex, Sec.
% sealed-noninterference and Sec. sealed-exercises [internal -- no
% script; the present mutation suite cannot exhibit this attack, which is
% the boundary's own point].
% Must distinguish: the honest path (gate reads committed s) from the
% laundering path (worker computes t=f(s) first); the shared prefix
% (load, read) and the shared honest gate (release: g(x)) from the one
% thing that differs (what x is bound to); that the checker's "submit"
% node is drawn identically in both lanes because it truly carries no
% payload in the model.
% Grammar chosen: paired paths on one grid (the before/after idiom, reused
% from the sibling fig-sealed-two-worlds two lanes below each other),
% divergence marked at the extra compute step the laundering lane detours
% through before rejoining the shared submit column.
% Grammar rejected: a single-lane flowchart with an if/else branch --
% that draws a choice the checker's automaton actually models and could
% in principle catch; the point is the reverse, that the checker's model
% has no branch here at all, only two lanes a human can compare on paper
% that the checker itself cannot.
% Five-second test: at the submit column, do the two lanes look
% different? They do not -- that is why the attack is priced, not caught.
% QA: beauty_lint B2 (arrow tip 0.3pt from the compute box) and B7
% (0.31pt sub-point centre offset between "release" and "g(x)") are the
% same two warnings the sibling fig-sealed-two-worlds.tex carries and
% leaves unfixed: an arrowhead meeting its node's edge, and a two-line
% node's bold/math baselines, both below print resolution. Justified,
% not fixed.
\begin{figure}[H]
\centering
\pdfigurehue{pdcobalt}%
\begin{tikzpicture}[pd figure,x=1cm,y=1cm]
  \def\lx{0}      % load column
  \def\rx{1.60}   % read column
  \def\cx{3.20}   % compute column -- laundering lane only
  \def\sx{4.85}   % submit column -- the checker's boundary
  \def\gx{6.65}   % release column
  \def\yH{3.35}   % honest lane
  \def\yL{0.40}   % laundering lane
  \node[pd title,anchor=east] at (-1.10,\yH) {honest path};
  \node[pd title,anchor=east] at (-1.10,\yL) {laundering path};
  % --- honest lane: load, read, submit s, release g(s) ----------------------
  \node[pd state,minimum width=1.45cm] (H1) at (\lx,\yH) {load};
  \node[pd state,minimum width=1.45cm] (H2) at (\rx,\yH) {read};
  \node[pd state,minimum width=1.45cm] (H3) at (\sx,\yH) {submit};
  \node[pd label,anchor=north] at (\sx,{\yH-.72}) {$x=s$};
  \node[pd focus state,minimum width=1.75cm,align=center] (H4) at (\gx,\yH) {release\\$g(x)$};
  \draw[pd focus arrow] (H1) -- (H2);
  \draw[pd focus arrow] (H2) -- (H3);
  \draw[pd focus arrow] (H3) -- (H4);
  % --- laundering lane: load, read, compute t=f(s), submit t, release g(t) -
  \node[pd state,minimum width=1.45cm] (L1) at (\lx,\yL) {load};
  \node[pd state,minimum width=1.45cm] (L2) at (\rx,\yL) {read};
  \node[pd breach state,minimum width=1.70cm,align=center] (L3) at (\cx,\yL) {compute\\$t=f(s)$};
  \node[pd state,minimum width=1.45cm] (L4) at (\sx,\yL) {submit};
  \node[pd label,anchor=south] at (\sx,{\yL+.72}) {$x=t$};
  \node[pd focus state,minimum width=1.75cm,align=center] (L5) at (\gx,\yL) {release\\$g(x)$};
  \draw[pd focus arrow] (L1) -- (L2);
  \draw[pd breach arrow] (L2) -- (L3);
  \draw[pd breach arrow] (L3) -- (L4);
  \draw[pd focus arrow] (L4) -- (L5);
  % --- the checker's boundary: the same "submit" node, no payload, either lane
  \begin{scope}[on background layer]
    \node[pd panel,fit={(H3)(L4)},inner xsep=9pt,inner ysep=11pt] (boundary) {};
  \end{scope}
  \node[pd title,anchor=south] at (boundary.north) {checker's boundary};
  \node[pd note,anchor=north west,align=left,text width=4.5cm] at (\cx-0.85,-0.95)
    {priced in the leakage budget, not caught by the checker};
\end{tikzpicture}
\caption{Two paths through the same honest gate, drawn as paired lanes. Above,
the honest path submits the committed secret $s$ directly, so the gate's
declared function releases $g(s)$. Below, the laundering path first computes
$t=f(s)$ for a worker-chosen $f$ and submits $t$ instead, so the same honest
gate releases $g(t)$ -- for a non-constant $f$ that maps an equal-parity pair
to a different-parity pair, this differs from what the honest path would have
released. The two lanes are drawn identically through \texttt{submit}, boxed
as the checker's boundary, because the model's \texttt{submit} action carries
no payload: the mutation suite that catches a leaky gate or a bypassed gate in
Theorem~\ref{thm:sealed-noninterference} has no variable to quantify over
here, so it cannot exhibit this attack at all. The residual is not eliminated;
it is priced as channel capacity in Section~\ref{sec:sealed-budget} [internal;
no script -- Exercise~\ref{ex:sealed-launder} states the model extension a
checker would need].}
\label{fig:sealed-laundering-fork}
\end{figure}
